\subsection{EMBEDDING REAL NUMBER SPACE IN PROJECTIVE SPACE}

Proposition 5 of no. 2 shows that if we are given a projective hyperplane H in \(P_{n}\) ( \(n \geqslant 0\) ), there is a homeomorphism of \(R^{n}\) onto the complement \(\complement H\) of this hyperplane. Once H has been chosen it is often convenient to identify \(R^{n}\) and \(\complement H\) by means of the homeomorphism defined in Proposition 5; the projective hyperplane H is then said to be ``at infinity'', and so are its points and subsets. Usually one takes H to be the ``coordinate'' hyperplane whose equation is \(x_{0}=0\) , and then the point \(z=(z_{i})\) of \(R^{n}\) is identified with the point of \(P_{n}\) whose homogeneous coordinates are \(1, z_{1}, z_{2}, \ldots, z_{n}\) .

Once this identification has been made, the closure in \(P_{n}\) of any affine linear variety V of p dimensions in \(R^{n}\) is a projective linear variety of p dimensions, not contained in the hyperplane at infinity, and identical with the projective linear variety generated by V. Conversely, every projective linear variety P of p dimensions which is not contained in the hyperplane at infinity has as its trace on \(R^{n}\) an affine linear variety of p dimensions, whose closure in \(P_{n}\) is P.

In the particular case \(n = 1\) , the hyperplane at infinity is a point; since \(\mathbf{P}_1\) is compact, it follows from Alexandroff's theorem (Chapter I, §9, no. 8, Theorem 4) that \(\mathbf{P}_1\) is homeomorphic to the space \(\tilde{\mathbf{R}}\) obtained by compactifying the locally compact space \(\mathbf{R}\) by the adjunction of a point (the ``point at infinity''). By Proposition 4 of §2, no. 4, we see also that the real projective line \(\mathbf{P}_1(\mathbf{R})\) is homeomorphic to the circle \(\mathbf{S}_1\) and to the torus \(\mathbf{T}\) .

On the other hand, if \(n > 1\) , \(\mathbf{P}_n(\mathbf{R})\) is not homeomorphic to \(\mathbf{S}_n\) , as we shall see later (cf.~Exercise 4).

The ``point at infinity'' of the space \(\tilde{R}\) is denoted by \(\infty\) , with no sign attached. It is important not to confuse \(\tilde{R}\) with the extended real line \(\overline{R}\) defined in Chapter IV, § 4, which has two ``points at infinity''; indeed \(\tilde{R}\) is homeomorphic to the quotient space of \(\overline{R}\) obtained by identifying the two points \(+\infty\) and \(-\infty\) .

